\documentclass[exercice,a4]{thedocument}%
\usepackage{texmaths-tikz}
\startdatakeys
\source{23GENMATMEAG1}
\cycle{4}
\competencesDuSocle{CA,RA,CO}
\connaissancesRequises{D2e*,D2d1}
\competencesTravaillees{C1a*,D2e*,D2d*}
\enddatakeys
\begin{document}%
Sur la figure ci-dessous~:
\begin{itemize}
    \item \mathspoint{BCDE} est un rectangle, \mathspoint{BAE} est un triangle
    rectangle en A~;
    \item la perpendiculaire à la droite \mathspoint{(CD)} passant par
    \mathspoint{A} coupe cette droite en \mathspoint{H}~;
    \item les droites \mathspoint{(AE)} et \mathspoint{(CD)} se coupent
    en F.
\end{itemize}
\begin{tikzpicture}
    \coordinate[label=below left:C](C)at(0,0);
    \coordinate[label=below:H](H)at(2.52,0);
    \coordinate[label=below:D](D)at(7,0);
    \coordinate[label=above left:B](B)at(0,4.2);
    \coordinate[label=below right:F](F)at(12.6,0);
    \coordinate[label=above:A](A)at(2.52,7.56);
    \path(A)--(F)coordinate[pos=4/9, label=above right:E](E);
    \draw(A)--(B)--(C)--(F)--cycle;
    \draw(B)--(E)--(D);
    \draw(A)--(H);
    \tkzMarkRightAngle[size=0.5,colored](B,A,E)
    \tkzMarkRightAngle[size=0.5,colored](C,B,E)
    \tkzMarkRightAngle[size=0.5,colored](D,C,B)
    \tkzMarkRightAngle[size=0.5,colored](E,D,C)
    \tkzMarkRightAngle[size=0.5,colored](B,E,D)
    \tkzMarkRightAngle[size=0.5,colored](D,H,A)
    \tkzMarkSegments[mark=s||,colored](A,B B,C)
    \tkzMarkSegments[mark=s|||,colored](B,E E,F)
\end{tikzpicture}\par\noindent
On donne~:
\begin{itemize}
    \item \mathspoint{AB}=\mathspoint{BC}= 4,2 cm~;
    \item \mathspoint{EB}=\mathspoint{EF}= 7 cm~;
\end{itemize}
\begin{enumerate}
    \item Montrer que l'aire du rectangle \mathspoint{BCDE} est égale à
    29,4 cm$^2$.
    \item \begin{enumerate}
        \item Montrer que la longueur \mathspoint{AE} est égale à 5,6 cm.
        \item Calculer l'aire du triangle \mathspoint{ABE}.
    \end{enumerate}
    \item\begin{enumerate}
        \item Montrer que les droites \mathspoint{(ED)} et
        \mathspoint{(HA)} sont parallèles.
        \item Calculer la longueur \mathspoint{AH}.
    \end{enumerate}
\end{enumerate}
\showthesource
\end{document}